\documentclass[landscape, 9pt, a4paper, oneside, twocolumn]{extarticle}
\usepackage{kotex}
\usepackage{amssymb}
\usepackage{amsmath}
\usepackage{import}
\usepackage{commath}
\usepackage{multicol}
\usepackage{multirow}

\usepackage{teamnote}

\teamnote{Petrozavodsk State University}{(Basalova, Gasova, Smirnov)}{PetrSU}

\DeclareMathOperator{\arccot}{arccot}
\newcommand{\floor}[1]{\left\lfloor #1 \right\rfloor}
\newcommand{\ceil}[1]{\left\lceil #1 \right\rceil}

\begin{document}

%\maketitlepage
\begin{minipage}{0.48\textwidth}
	\begin{multicols}{2}
		\tableofcontents
	\end{multicols}
\end{minipage}

% Make Pagebreak if you want.
%\pagebreak 

%\tableofcontents

\newpage

\ShowUsage
\ShowComplexity
\HideAuthor


\newpage

\section{Общее}

\Algorithm
{Шаблон}
{}
{}
{cpp}{source/General/template.cpp}
{}

\Algorithm
{Строка компиляции}
{}
{}
{cpp}{source/General/debug_compilation.txt}
{}

\Algorithm
{Рандом}
{}
{}
{cpp}{source/General/random.cpp}
{}

\Algorithm
{Built-in GCC Stuff}
{}
{}
{cpp}{source/General/gcc-stuff.txt}
{}

\Algorithm
{\_\_int128}
{}
{}
{cpp}{source/General/int128.cpp}
{}

\Algorithm
{PBDS: Ordered set}
{}
{}
{cpp}{source/General/PBDS.cpp}
{}

\Algorithm
{Стресс-тестирование}
{}
{}
{py}{source/General/stress.py}
{}

\section {Графы}

\Algorithm
{Поиск мостов}
{}
{}
{cpp}{source/Graphs/bridges.cpp}
{}

\Algorithm
{Поиск точек сочленения}
{}
{}
{cpp}{source/Graphs/cutpoints.cpp}
{}

\Algorithm
{LCA: sparse table}
{}
{}
{cpp}{source/Graphs/lca.cpp}
{}

\Algorithm
{MST: алгоритм Прима}
{}
{}
{cpp}{source/Graphs/prim.cpp}
{}

\Algorithm
{Алгоритм Куна}
{}
{}
{cpp}{source/Graphs/kuhn.cpp}
{}

\Algorithm
{Эйлеров обход}
{}
{}
{cpp}{source/Graphs/euler.cpp}
{}

\Algorithm
{2-SAT}
{}
{}
{cpp}{source/Graphs/sat.cpp}
{}

\Algorithm
{Maxflow: алгоритм Диница}
{}
{}
{cpp}{source/Graphs/dinic.cpp}
{}

\Algorithm
{Minimal cut}
{}
{}
{cpp}{source/Graphs/mincut.cpp}
{}

\Algorithm
{Mincost flow}
{}
{}
{cpp}{source/Graphs/mincost.cpp}
{}

\section{Математика}

\input source/Mathematics/formulas.tex

\Algorithm
{FFT}
{}
{}
{cpp}{source/Mathematics/FFT.cpp}
{}

\Algorithm
{Нахождение определителя}
{}
{}
{cpp}{source/Mathematics/determinant.cpp}
{}

\Algorithm
{Метод Гаусса}
{}
{}
{cpp}{source/Mathematics/gauss.cpp}
{}

\Algorithm
{Китайская теорема об остатках}
{}
{}
{cpp}{source/Mathematics/CRT.cpp}
{}

\Algorithm
{Тест Миллера-Рабина}
{}
{}
{cpp}{source/Mathematics/miller-rabin.cpp}
{}

\Algorithm
{Функция Эйлера}
{}
{}
{cpp}{source/Mathematics/euler_fun.cpp}
{}

\Algorithm
{Линейное решето Эратосфена}
{}
{}
{cpp}{source/Mathematics/linear-sieve.cpp}
{}

\Algorithm
{Диофантовы уравнения}
{}
{}
{cpp}{source/Mathematics/diofant.cpp}
{}

\Algorithm
{Первообразный корень}
{}
{}
{cpp}{source/Mathematics/generator.cpp}
{}

\subsection {Теория игр}
\input source/GameTheory/sprague-grandi.tex

\section{Оптимизации динамики}
\Algorithm
{LineContainer (aka $\mathcal{O} (N \log{N})$ CHT)}
{}
{}
{cpp}{source/dp/line_container.cpp}
{}

\Algorithm
{Convex Hull Trick}
{}
{}
{cpp}{source/dp/cht.cpp}
{}

\Algorithm
{Оптимизация Кнута}
{}
{}
{cpp}{source/dp/knut.cpp}
{}

\section{Структуры данных}
\Algorithm
{Декартово дерево с прибавлением на отрезке}
{}
{}
{cpp}{source/DataStructures/cartesian.cpp}
{}

\Algorithm
{Дерево Фенвика}
{}
{}
{cpp}{source/DataStructures/fenwick.cpp}
{}

\Algorithm
{2D дерево Фенвика}
{}
{}
{cpp}{source/DataStructures/2d-fenwick.cpp}
{}

\Algorithm
{Персистентное ДО}
{}
{}
{cpp}{source/DataStructures/pers-segtree.cpp}
{}

\section{Строки}

\Algorithm
{Z-функция}
{}
{}
{cpp}{source/Strings/z.cpp}
{}

\Algorithm
{Префикс-функция}
{}
{}
{cpp}{source/Strings/prefix.cpp}
{}

\Algorithm
{Бор}
{}
{}
{cpp}{source/Strings/trie.cpp}
{}

\Algorithm
{Суффиксный массив}
{}
{}
{cpp}{source/Strings/suf_array.cpp}
{}

\Algorithm
{Суффиксный автомат}
{}
{}
{cpp}{source/Strings/suf_automata.cpp}
{}

\section{Геометрия}

\Algorithm
{Шаблон}
{}
{}
{cpp}{source/Geometry/geometry.cpp}
{}

\Algorithm
{Выпуклая оболочка}
{}
{}
{cpp}{source/Geometry/convex.cpp}
{}

\Algorithm
{Полезные функции}
{}
{}
{cpp}{source/Geometry/functions.cpp}
{}

\end{document}
